% RESEARCH_DRAFT_SHA256: be58dce8513d26cf7b4da16257e95fe174b19c6e946c126b72720922e3c67616
% Layout and source-reference wrappers expanded for the web renderer.
\documentclass[11pt]{article}
\usepackage[margin=25mm]{geometry}
\usepackage{amsthm}
\usepackage{fontspec,unicode-math}
\setmainfont{Latin Modern Roman}
\setmathfont{Latin Modern Math}
\usepackage{amsmath,amssymb,mathtools}
\usepackage{xurl,hyperref}
\hypersetup{colorlinks=true,urlcolor=blue,linkcolor=blue}
\setcounter{secnumdepth}{0}
\setlength{\parindent}{0pt}
\setlength{\parskip}{5pt}
\setlength{\emergencystretch}{3em}
\newcommand{\sref}[2]{\hyperlink{source:#1:#2}{[S#2]}}
\newcommand{\cataloguescope}[1]{\paragraph{Source question.}#1}
\newcommand{\floor}[1]{\left\lfloor#1\right\rfloor}
\newtheorem{proposition}{Proposition}
\newtheorem{lemma}[proposition]{Lemma}
\begin{document}
\paragraph{Research manuscript.} OpenAI \texttt{gpt-6-astra}, reasoning \texttt{ultra}; 11 October 2026.
\paragraph{Preserved source body.} The complete catalogue mathematical text and sources follow with layout and source links expanded, including source qualifications and its historical preparation notes. New research begins at the section headed Research attempt.

% UnsolvedMath ID 20001213; source catalogue code AIM-COMPUTATION-0051
\section{A gap-move bound and the exact decomposable-model spectrum}\label{20001213}
{\small UnsolvedMath ID 20001213 · Computer Science}
\subsection{Definitions and mathematical statement}

\paragraph{Shared notation.}Standard number systems are denoted by $\mathbb N,\mathbb Z,\mathbb Q,\mathbb R,\mathbb C$; index ranges, computational models and specialized notation follow the cited problem source and the definitions below.


For fixed margins $b$, consider the nonnegative integral tables in the fiber $\mathcal F_b=\{x:Ax=b\}$. For each cell $i$, its attainable-value set is $\{x_i:x\in\mathcal F_b\}$. The workshop's spectrum records these possible cell-value sets as the margins vary for the given hierarchical model.
\paragraph{Statement recorded in the supplied source.}
Problem 37. Find the spectra of various models and classes of models. \href{https://aimath.org/WWN/compalgstat/compalgstat.pdf}{[S2]}
\paragraph{Scope qualification.}
This is a classification question across the stated model families; the spectrum is not the set of eigenvalues of a matrix.
\paragraph{Source status.}
The supplied catalogue labels this record partially solved. Its source-status metadata was checked against the supplied snapshot on 2026-10-11. This import does not independently establish current open status or mathematical correctness.
\subsection{Short English statement}
Which sets of integer cell values can occur in tables with fixed marginal totals?
\subsection{Sources}
\begin{itemize}
\item[S1] Ulam AI and the UnsolvedMath Contributors, UnsolvedMath: A Curated Collection of Open Mathematics Problems, record 20001213 (AIM-COMPUTATION-0051). Catalogue curation and metadata: CC BY 4.0. The supplied source question is retained; definitions and the English summary are editorial additions.
\url{https://huggingface.co/datasets/ulamai/UnsolvedMath}.
\item[S2] AIM Workshop Problem Lists, original problem source linked by catalogue record AIM-COMPUTATION-0051
\url{https://aimath.org/WWN/compalgstat/compalgstat.pdf}.
\end{itemize}
\paragraph{Preparation and review.}
Original question attributed to the cited source and UnsolvedMath contributors. Codex (GPT-6; exact serving revision not exposed) assisted with definitions, English summaries, source inspection and formatting. No new human mathematical authorship or independent proof verification is asserted. Canonical ID 20001213 and source code AIM-COMPUTATION-0051 checked against the complete supplied registry and existing repository statements; known flagged repetitions and obvious identical conjectures excluded. Full mathematical-equivalence/literature review has not been performed. Statement preparation and final review followed contribution prompts C and D; automated checks verify structure and identity, not open status.
\label{end:20001213}


\section{Research attempt}
\subsection{Literature and scope}
The source uses ``spectrum'' for possible cell counts in a nonnegative integer fiber. I tested whether high Markov degree itself forces gaps in a cell spectrum. Binary tables with all proper margins fixed give an exact counterexample to that inference. This is a standard parity-family calculation, not a classification of the source's various models.
\begin{proposition}
For $r\ge2$ binary variables with every $(r-1)$-variable margin fixed, every cell spectrum is an integer interval although the required Markov degree is $2^{r-1}$.
\end{proposition}
\begin{proof}
For two tables in the same fiber, their difference $v$ satisfies $v_x+v_{x\oplus e_j}=0$ for every bit $j$. Therefore $v_x=t(-1)^{\sum_jx_j}$ for the integer $t=v_0$. Write $u_x=(-1)^{\sum_jx_j}$ and choose a feasible table $z$. Conversely every integer multiple of $u$ preserves these margins. Hence the fiber is exactly
\[
\{z+tu:t\in\mathbb Z,\ -m_+\le t\le m_-\},\qquad
m_+=\min_{u_x=1}z_x,\quad m_-=\min_{u_x=-1}z_x.
\]
For a specified cell its value is $z_x+t$ or $z_x-t$, so every integer between its endpoints occurs. The primitive move $u$ connects the whole fiber and has degree $2^{r-1}$. In the fiber through $u^+$, the only two tables are $u^+,u^-$, because the allowed parameter values are $-1,0$. Connecting this fiber forces a move of that degree, proving the claimed necessity.
\end{proof}
\subsection{Failed route and original-scope comparison}
The growth of total positive mass in a move does not force a large change in an individual coordinate: every coordinate of $u$ is $\pm1$. Thus any argument transferring large Markov degree directly into cell gaps fails here. For $r=4$ and the constant table $z_x=3$, the allowed parameter values are $-3,-2,\ldots,3$ and every cell spectrum is $\{0,\ldots,6\}$, while degree eight is required across all fibers of the model.
\subsection{Verification and remaining gaps}
The example's interval and degree were checked by the displayed formulas. The proof classifies spectra only for this rank-one-kernel family and does not extend to arbitrary hierarchical models, structural zeros or more general prescribed statistics. No new original-target progress is asserted.

\section{References and source attribution}
The original question is credited to the AIM workshop list prepared by Seth Sullivant (2004), \url{https://aimath.org/WWN/compalgstat/compalgstat.pdf}. The source question and catalogue attribution are reproduced above. This list was accessed on 11 October 2026; it is historical evidence of the question, not a certification of present open status.


\paragraph{Integration review.} Codex (GPT-6; exact serving revision not exposed) reviewed the written arguments and expanded layout and source-reference wrappers. Repository metadata and the production web renderer were checked using Python 3.13 and JavaScript with MathJax 3.2.2. This model-assisted review is not independent human proof verification; no PDF compilation or formal proof checking is asserted.

\section{Completion Estimate}
\noindent\textbf{COMPLETION: 0\%}
No substantiated new progress toward the entire original target is claimed. The worked results and failed route are delimited above; a known result, illustrative model, or auxiliary obstruction is not counted as original-target completion.

\section{Final status and attribution}
\textbf{Attempt status: unresolved.} This describes this research attempt, not a new ruling on the source catalogue's status or on external literature claims.
The mathematics and this self-audit were generated by OpenAI \texttt{gpt-6-astra}, reasoning \texttt{ultra}, on 11 October 2026. Codex (GPT-6; exact serving revision not exposed) orchestrated the batch and contributed formatting/repository review. No human mathematical authorship, independent human proof checking, or formal verification is claimed. The historical preparation paragraphs in the reproduced catalogue body concern that body only.

Contribution prompt A was applied to the complete definition and supplied source record, followed by the adversarial checks of prompt B. Local reproduction files are \texttt{batch005/checks.py}, \texttt{computations.json}, and the per-ID audit JSON. They were not public links at preparation time. Numerical checks are finite and do not prove unrestricted results. No LaTeX engine was found locally; compilation was not performed by this research agent.

\end{document}
